\section{Introdução}
A orientação da Entrega XIII solicita que a equipe elabore um plano de crescimento e escalabilidade adequado ao estágio atual da startup, contemplando diagnóstico, duas estratégias de crescimento, ações de escalabilidade e plano de ação com recursos e indicadores \cite{pucxiii2026}. Esta entrega consolida a Entrega XII, as versões Alfa da XIII e os artefatos anteriores do repositório do projeto. Seu objetivo não é apresentar crescimento como certeza, mas transformar hipóteses do modelo de negócio em experimentos verificáveis.

A Entrega XII estruturou o pitch em onze lâminas, com narrativa problema--solução--mercado--negócio--evidência--futuro. A continuidade é preservada: a XIII aprofunda o caminho operacional e econômico que dá sustentação ao futuro anunciado no pitch. A documentação da XII registra que as evidências de protótipo e testes são exploratórias e que projeções devem ser identificadas como hipóteses, princípio mantido nesta revisão.

A literatura distingue crescimento, entendido como aumento de usuários ou receita, de escalabilidade, que requer absorver esse aumento sem elevação proporcional de custos \cite{sanasi2023,platformscaling2023}. Para plataformas bilaterais, a expansão depende também da coordenação dos grupos conectados e da criação de valor entre eles \cite{rochet2003,eisenmann2006}. Assim, a pergunta que orienta o plano é: Como validar uma praça local e elevar a contribuição econômica por usuário antes de replicar a operação?